% 385_Higgs_Vakuum_Weg_En_ch.tex
% Johann Pascher, 1 Oct 2026
% The Higgs-vacuum route to xi: origin, formulas, status

\providecommand{\BM}{\textbf{[B]}}
\providecommand{\KM}{\textbf{[K]}}
\providecommand{\SM}{\textbf{[S]}}
\providecommand{\XM}{\textbf{[X]}}
\providecommand{\QM}{\textbf{[Q]}}

\chapter*{The Higgs--vacuum route to \texorpdfstring{$\xi$}{xi}}
\addcontentsline{toc}{chapter}{The Higgs-vacuum route to xi}

\begin{abstract}
The first route to the relations of FFGFT ran through the vacuum and the Higgs field. In spring
2025 it produced a formula linking the Higgs quantities to the electric constant
$\varepsilon_0$: $\lambda_h^2v^2e^2/(64\pi^4\varepsilon_0\hbar c\,m_h^2)$. This document traces the
formula through the version history from February 2025 to today, recomputes every stage and
states the current status. Written out, the vacuum formula is exactly $\xi_{\text{EFT}}\cdot\alpha$
with $\xi_{\text{EFT}}=\lambda_h^2v^2/(16\pi^3m_h^2)$; with the FFGFT convention $\alpha=1$ it
becomes today's Higgs route. With the self-coupling $\lambda_h=m_h^2/(2v^2)$ one obtains
$\xi_{\text{EFT}}=m_h^2/(64\pi^3v^2)=1.30\times10^{-4}$, about $2.3\,\%$ below $\xi=4/30000$. The
Higgs geometry is therefore a consistency check of $\xi$ to within a few per cent \KM, not an exact
derivation. Not sustainable are the early equating of the expression with~1, a short-lived variant
with $\mu_0/\varepsilon_0$, and the later shortcut to $1/(16\pi^3)$ \XM. The factor $64\pi^4$ of the
early HTML versions comes from the $4\pi$ in $4\pi\varepsilon_0$.
\end{abstract}

% ============================================================
\section*{Status markers}
% ============================================================
\begin{center}
\begin{tabular}{@{}lp{11cm}@{}}
\textbf{[B]} & Algebraically proven. \\[2pt]
\textbf{[K]} & Numerically confirmed (check script, comparison with measured values). \\[2pt]
\textbf{[S]} & Assumption, open or speculative. \\[2pt]
\textbf{[X]} & Refuted or not sustainable. \\[2pt]
\textbf{[Q]} & Convention. \\
\end{tabular}
\end{center}
Reference values: CODATA 2022, PDG 2024 ($m_h=125.20$\,GeV, $v=246.22$\,GeV). The sources of the
early formulas are commits of the repository; the files now exist only in the version history.
Check script: \texttt{pruef\_385\_higgs\_vakuum.py}.

% ============================================================
\section{Starting idea: vacuum, Higgs field and the speed of light}
% ============================================================

In February 2025 the starting point was a qualitative thesis: the vacuum is not empty, and the
Higgs field together with the background radiation determines the maximal speed. The link ran
through the vacuum constants,
\begin{equation}
  c=\frac{1}{\sqrt{\mu_0\varepsilon_0}},\qquad Z_0=\sqrt{\frac{\mu_0}{\varepsilon_0}}=376.73\ \Omega ,
\end{equation}
still without a formula for a parameter (first upload of 27~March 2025, document on the
Schwarzschild metric). In parallel, notes appeared expressing $\mu_0$ and $\varepsilon_0$ through
$\alpha$, $e$ and $c$, $\alpha=e^2/(4\pi\varepsilon_0\hbar c)$. These two lines, Higgs and vacuum
on one side, vacuum constants and $\alpha$ on the other, are the origin of the route traced here.

% ============================================================
\section{The chain of formulas in 2025}
% ============================================================

\begin{center}
\small
\begin{tabular}{@{}p{2.2cm}p{6.4cm}p{5.4cm}@{}}
\toprule
\textbf{Date} & \textbf{Formula} & \textbf{Finding} \\
\midrule
31 Mar 2025 & $\beta=\dfrac{\lambda_h^2v^2}{16\pi^3c^3}\cdot\dfrac{\hbar}{m_h^2}\cdot\dfrac1{r_0}$ & factor $16\pi^3$ appears; not dimensionless in SI \\[8pt]
5 Apr 2025 & $\xi=\dfrac{\lambda_h^2v^2}{16\pi^3m_h^2}\approx1.33\times10^{-4}$ & correctly evaluated $1.30\times10^{-4}$ \\[8pt]
6 Apr 2025 & $\dfrac{\lambda_h^2v^2e^2}{64\pi^4\varepsilon_0\hbar c\,m_h^2}=1$ & $=\xi_{\text{EFT}}\,\alpha$, not 1 \\[8pt]
18 Apr 2025 & $\beta_T=\dfrac{\lambda_h^2v^2}{4\pi^2\lambda_0^2\alpha_0}$, dimension $(\mu_0/\varepsilon_0)^k$ & dimensionally inconsistent \\[8pt]
25 May 2025 & $\varepsilon_0$ eliminated via $e^2=4\pi\varepsilon_0$: $64\pi^4\to16\pi^3$ & return to the Higgs expression \\[8pt]
from 1 Jun 2025 & $\xi=\dfrac{\lambda_h^2v^2}{64\pi^4m_h^2}=1.04\times10^{-5}$ & $4\pi$ from $4\pi\varepsilon_0$ left over \\[8pt]
22 Aug 2026 & $\lambda_h=m_h/v\Rightarrow\xi=1/(16\pi^3)$, iteration to $4/30000$ & shortcut, not sustainable \\
\bottomrule
\end{tabular}
\end{center}

\subsection{The Higgs expression}

On 31~March 2025 an expression with the Higgs quantities and the factor $16\pi^3$ appears for the
first time, then as a parameter $\beta$ with a radius $r_0$. On 5~April 2025, with
$r_0=\xi\,\ell_P$, natural units and $\beta_T=1$, it becomes
\begin{equation}
  \xi=\frac{\lambda_h^2v^2}{16\pi^3m_h^2}.
\end{equation}
The document gives $1.33\times10^{-4}$; computed from the inputs it is $1.30\times10^{-4}$. The same
document already contains the rewriting with $m_h^2=2\lambda_hv^2$: $\xi=\lambda_h/(32\pi^3)$. The
expression is dimensionless and correctly formed \BM.

\subsection{The vacuum formula}

On 6~April 2025 the document connects, via natural units with $\alpha=1$, the Higgs expression
with the fine-structure constant in its SI form:
\begin{equation}
  \beta_T\,\alpha_{\text{EM}}\approx\frac{\xi\,\lambda_h^2v^2}{16\pi^3m_h^2}\cdot\frac1\xi\cdot
  \frac{e^2}{4\pi\varepsilon_0\hbar c}
  =\frac{\lambda_h^2v^2e^2}{64\pi^4\varepsilon_0\hbar c\,m_h^2}=1 .
  \label{eq:vakuum}
\end{equation}
The same document contains $c=1/\sqrt{\mu_0\varepsilon_0}$, $Z_0=\sqrt{\mu_0/\varepsilon_0}$ and
$h=2\pi\hbar\mu_0c^2$. This is the formula in which the Higgs sector and the vacuum constants
appear together for the first time.

\subsection{The variant with \texorpdfstring{$\mu_0/\varepsilon_0$}{mu0/eps0}}

On 18~April 2025 an English working draft attempts an ``electromagnetic derivation'' of
$\beta_T$ from $\mu_0$ and $\varepsilon_0$:
$\beta_T^{\text{nat}}=\lambda_h^2v^2/(4\pi^2\lambda_0^2\alpha_0)$ with the Higgs Compton length
$\lambda_0=\hbar/(m_hc)$ and a dimension factor $(\mu_0/\varepsilon_0)^k$, $k=0$. The file was
deleted two days later and never adopted.

\subsection{Return and offshoot}

On 25~May 2025 $\varepsilon_0$ is eliminated via $e^2=4\pi\varepsilon_0$ ($\alpha=1$,
$\hbar=c=1$); $64\pi^4$ turns back into $16\pi^3$, and the expression is again the Higgs
expression of 5~April. From 1~June 2025, however, the HTML versions and Doc.~049 contain
$\xi=\lambda_h^2v^2/(64\pi^4m_h^2)=1.04\times10^{-5}$: here $\varepsilon_0$ was dropped, but the
$4\pi$ from $4\pi\varepsilon_0$ was not cancelled.

\subsection{The shortcut of 2026}

When Doc.~320 was created on 22~August 2026, $\lambda_h=m_h/v$ was inserted. Everything then
reduces to $1/(16\pi^3)$, and an iteration is supposed to lead from there to the fixed point
$4/30000$; A130 and Doc.~133 are cited as sources. This substitution confuses the Higgs
self-coupling with a Yukawa coupling $y=m/v$; the same identification $y=\lambda_h$ already appears
in Doc.~097 in 2025. A130 contains the Higgs route in no version, Doc.~133 contains no iteration.

% ============================================================
\section{What the vacuum formula means}
% ============================================================

Written out, the right-hand side of \eqref{eq:vakuum} is a product of two dimensionless
factors:
\begin{equation}
  \frac{\lambda_h^2v^2e^2}{64\pi^4\varepsilon_0\hbar c\,m_h^2}
  =\underbrace{\frac{\lambda_h^2v^2}{16\pi^3m_h^2}}_{\xi_{\text{EFT}}}\cdot
   \underbrace{\frac{e^2}{4\pi\varepsilon_0\hbar c}}_{\alpha}
  =\xi_{\text{EFT}}\cdot\alpha\quad\BM .
\end{equation}
The formula is dimensionally correct. The vacuum constant $\varepsilon_0$ enters only through
$\alpha$, and the factor $64\pi^4=16\pi^3\cdot4\pi$ arises because the $4\pi$ from
$4\pi\varepsilon_0$ moves into the denominator. In SI the expression gives
$\xi_{\text{EFT}}\,\alpha=9.5\times10^{-7}$ \KM. With the FFGFT convention $\alpha=1$, i.e.
$e^2=4\pi\varepsilon_0\hbar c$, it is exactly $\xi_{\text{EFT}}$ \BM. The Higgs--vacuum route and
today's Higgs route are therefore the same calculation: the vacuum constants disappear in natural
units, and what remains is the Higgs expression that checks $\xi$. Doc.~310 reads it the same way:
in natural units $\mu_0$ and $\varepsilon_0$ disappear, and only $\xi$ remains.

Not sustainable is the equating with~1 in \eqref{eq:vakuum} \XM. In the first step
$\xi\cdot(\ldots)\cdot(1/\xi)$ is cancelled, while $\beta_T$ itself contains the factor $1/\xi$.
Read correctly, $\beta_T=\xi_{\text{EFT}}/\xi=0.977$, close to but not exactly~1, and
$\beta_T\alpha=0.0071$.

The variant with $\mu_0/\varepsilon_0$ of 18~April is dimensionally inconsistent \XM:
$\lambda_0^{-2}$ has the dimension energy$^2$, so $v^2/\lambda_0^2$ has energy$^4$; with
$\lambda_0=1/m_h$ the expression gives about $4\times10^5$\,GeV$^4$, neither~1 nor $\xi$. The
factor $(\mu_0/\varepsilon_0)^k=Z_0^{2k}$ appears only in the dimension line and vanishes for
$k=0$.

% ============================================================
\section{Current status}
% ============================================================

With the Standard Model self-coupling $\lambda_h=m_h^2/(2v^2)=0.129$, the Higgs expression becomes
\begin{equation}
  \xi_{\text{EFT}}=\frac{\lambda_h^2v^2}{16\pi^3m_h^2}=\frac{m_h^2}{64\pi^3v^2}
  =\frac{\lambda_h}{32\pi^3}=1.30\times10^{-4}\quad\KM ,
\end{equation}
about $2.3\,\%$ below $\xi=4/30000=1.333\times10^{-4}$. It depends only on the ratio $m_h/v$.
Depending on the inputs the gap lies between $2.2$ and $2.4\,\%$; the frequently quoted $2.6\,\%$
arises only with $\lambda_h$ rounded to $0.129$:
\begin{center}
\small
\begin{tabular}{@{}lcc@{}}
\toprule
\textbf{Inputs} & $\xi_{\text{EFT}}$ & \textbf{Gap to }$4/30000$ \\
\midrule
$m_h=125.20$, $v=246.22$\,GeV (PDG 2024) & $1.303\times10^{-4}$ & $-2.3\,\%$ \\
$m_h=125.25$, $v=246.22$\,GeV & $1.304\times10^{-4}$ & $-2.2\,\%$ \\
$m_h=125.1$, $v=246.2$\,GeV & $1.301\times10^{-4}$ & $-2.4\,\%$ \\
$\lambda_h=0.129$ rounded, $m_h=125.1$ & $1.299\times10^{-4}$ & $-2.6\,\%$ \\
\bottomrule
\end{tabular}
\end{center}
The Higgs geometry is therefore a consistency check of $\xi$ to within a few per cent \KM\
(Doc.~354, A190, 310). The deviation is interpreted structurally \SM\ (R59): the SI values of
$m_h$ and $v$ are scheme dependent, and fractal corrections cannot be cleanly separated from the
loop structure of the Standard Model. The Higgs route is not an exact derivation of $4/30000$. It
is, however, an independent check: it uses only the Higgs mass and the electroweak scale, no lepton
mass and no Galois counting.

% ============================================================
\section{What holds and what does not}
% ============================================================

\begin{center}
\small
\begin{tabular}{@{}p{7.2cm}p{1.2cm}p{5.4cm}@{}}
\toprule
\textbf{Statement} & \textbf{Status} & \textbf{Reason} \\
\midrule
Vacuum formula $=\xi_{\text{EFT}}\cdot\alpha$ & \BM & written out; $\alpha=e^2/(4\pi\varepsilon_0\hbar c)$ \\
with $\alpha=1$ equal to the Higgs expression & \BM & $e^2=4\pi\varepsilon_0\hbar c$ \\
$\xi_{\text{EFT}}=m_h^2/(64\pi^3v^2)=1.30\times10^{-4}$, about $2.3\,\%$ below $\xi$ & \KM & consistency check \\
Deviation structural (scheme, fractal correction) & \SM & R59 \\
Vacuum formula $=1$ & \XM & cancellation of $1/\xi$ inconsistent \\
Variant with $\mu_0/\varepsilon_0$ & \XM & dimension energy$^4$ \\
$\xi=\lambda_h^2v^2/(64\pi^4m_h^2)=1.04\times10^{-5}$ & \XM & $4\pi$ from $4\pi\varepsilon_0$ not cancelled \\
$\lambda_h=m_h/v$, $\xi=1/(16\pi^3)$, iteration to $4/30000$ & \XM & 15 times $\xi$; iteration runs to 0 \\
Higgs route as exact derivation of $4/30000$ & \XM & gap 2--3\,\% \\
\bottomrule
\end{tabular}
\end{center}

% ============================================================
\section{Conclusion}
% ============================================================

The Higgs--vacuum route that stood at the beginning of FFGFT holds. Its vacuum formula of April
2025 is, written out correctly, the product of the Higgs expression and the fine-structure
constant; with $\alpha=1$ it is today's Higgs route, and this route meets $\xi$ to about
$2.3\,\%$. The errors of the line of development lie in side steps: the equating with~1, a
dimensionally wrong variant, an uncancelled $4\pi$ and a late shortcut. As a derivation of the
exact number $4/30000$ the route is not sufficient; as an independent check of $\xi$ from the
Higgs sector it stands.

\vspace{1em}
\noindent\textit{Check script:} \texttt{2/python/Dok385\_Skripte/pruef\_385\_higgs\_vakuum.py}
--- 17/17.
